%!TEX program = lualatex
\documentclass[ngerman,headheight=70pt]{scrartcl}
\usepackage[ngerman]{babel}
\usepackage{blindtext}
\usepackage[utf8]{luainputenc}
\usepackage{scrlayer-scrpage}
\usepackage[T1]{fontenc}
\usepackage{fontspec}
\usepackage{multicol}
\usepackage{enumitem}
\usepackage[german=quotes]{csquotes}
\usepackage[bottom=2cm,footskip=8mm]{geometry}
\setmainfont{Source Sans Pro}
\MakeOuterQuote{"}

\setlength{\parindent}{0pt}
\setlength{\parskip}{6pt}

\pagestyle{plain.scrheadings}
\thispagestyle{scrheadings}
\lohead{\vspace*{-1.5cm}\includegraphics[scale=0.4]{../images/stupa_logo_rot.pdf}}
\rohead{\\Präsidium des\\Studierendenparlaments}

\lofoot[Universität Hamburg · Präsidium des Studierendenparlaments\\\hrule Von-Melle-Park 5 · D-20146 Hamburg]{Universität Hamburg · Präsidium des Studierendenparlaments\\\hrule Von-Melle-Park 5 · D-20146 Hamburg}
\cofoot[]{}
\cohead[\thepage]{}

\renewcommand{\thesubsection}{TOP \arabic{subsection}}
\renewcommand{\thesubsubsection}{\arabic{subsubsection}.}
\setcounter{subsection}{-1}

\begin{document}
    UHH · StuPa-Präsidium · Von-Melle-Park 5 · D-20146 Hamburg

    %%% PAGE ONE (EDIT LATER) %%%%
    % Hier den Titel des Protokolls eintragen
    \section*{Protokoll der 12. Sitzung des Studierendenparlaments vom 15. Dezember 2016}

    % Protokollanten eintragen
    \textbf{Protokoll: Ramon Weilinger}\\
    \textbf{Ort: VMP5 0079}\\
    \textbf{Beginn: 20.47 Uhr}\\
    \textbf{Ende: 21.30 Uhr}

    % Wohl die komplizierteste Stelle im Protokoll
    % Als Faustregel sollten nicht mehr als 3 Namen auf einer Zeile stehen
    % Einfach an dieser Vorlage orientieren
    % Im Idealfall ist hier nicht viel zu ändern, da kaum jemand fehlt.
    % Rücktritte müssen natürlich pro Protokoll angepasst werden
    \vspace{0.5cm}
    \begin{tabular}{ll}
        Anwesend: & \\
            RCDS (5 Sitze): & Ramon Weilinger \\
             CampusGrün (14 Sitze): & Laura Franzen, Geoffrey Youett, Philipp Droll, \\
                                   & Yasemin Günther, Christian Lagod, Jim Martens, \\
                                   & Svenja Horn, Martin Sievert, Armin Günther \\
             BGZM-Liste (2 Sitze) : & \\
             WiWi (2 Sitze): & \\
             Unicorns (5 Sitze): &    \\
             Liste LINKS (3 Sitze): & Gunhild Berdal, Till Petersen, Sinah Mielich \\
             HWP (2 Sitze): &  \\
             MIN (4 Sitze): &  \\
             SDS* (3 Sitze): & Mena Winkler, Jacob Petersein, Artur Brückmann \\
             Bart-LISTE (2 Sitze): & Timo Zeimet, Malte Peters \\
             LHG (1 Sitz): & Tobias Heisig \\
             harte zeiten (1 Sitz): & Tobias Berking \\
             Jura (1 Sitz): & \\
             AL (2 Sitze): &  \\
            & \\
        Entschuldigt: & Antonia Niecke, Ramin Shakiba, Jennifer Maack, \\
                      & Benjamin Welling (RCDS), Mario Moldenhauer, Martin Sievert (CG) \\
                      & Jakob Pape, Karen Martirosian (BGZM), Claas-Friso Hente, \\
                      & Charleen Lorenz (WiWi), Katharina Kucza, Johannes Peplow, \\
                      & Annkathrin Löffler, Andreas Hartkamp (Unicorns), Hakan Tasdemir, \\
                      & Gökce Celik (HWP), Ailina Salten, Lotte Rullkötter, \\
                      & Nabil Jabareen, Jan Detampel (MIN), Karima Schulze\\
                      & Henri Weber (AL) \\
                      &\\
        Unentschuldigt abwesend: & Elena Rysikova, Melf Johannsen, Ines Friedrich \\
                                & Dean Haltmeyer (CG), Marielle Hermstrüwer (Unicorns), \\
                                & Maurice Martens (Jura)\\
        Rücktritte: & Daniel Runge \rightarrow Malte Peters (Bart-Liste) \\
        Mandatsverluste: & \\
    \end{tabular}
    \newpage
    %%% Vorgeschlagene Tagesordnung %%%
    \underline{Vorgeschlagene Tagesordnung}
    \begin{enumerate}[label={\textbf{Top \theenumi}},leftmargin=*]
        \item Beitragsordnung (V1617-046) (30 Min.)
        \item VG-Wort (V1617-047) (30 Min.)
        \item Verschiedenes (5 Min.)
    \end{enumerate}

    \newpage

    % Hier beginnt Teil A. Bei Fortsetzungen kommt hier nur der Teil über
    % die Feststellung der endgültigen Fassung der Tagesordnung.
    % Für sonstige Sitzungen ist dies bereits das Gerüst für die Formalia.
    % Man muss nur noch die protokollierten Inhalte eintragen.

    {\Large \underline{Teil A}}

    \subsection{Formalia}

    \subsubsection{Geschäftsbericht Präsidium}

    Gunhild eröffnet die Sitzung um 20.47 Uhr und berichtet für das Präsidium:

    Das Präsidium befasst sich aktuell vor allem mit der Vorbereitung und Durchführung der StuPa-Wahlen.
    \begin{itemize}
        \item Am kommenden Wochenende werden insbesondere die Wohnheime mit
        Wahlzeitungen ausgestattet. Leider gab es abermals ein Problem beim
        Druck, weshalb die Wahlzeitung später geliefert wurde als geplant.
        \item Durch Überlastung in der Universitätsverwaltung sei auch die
        STiNE-Mail verspätet verschickt worden.
    	\item Das Präsidium hat sich mit der Universitätsverwaltung über die
        Foyernutzung zur Urnenwahl verständigt. Aufgrund der
        Brandschutzbestimmungen wurde der Vorschlag erarbeitet, dass Tische und
        Stühle aus Aluminium für die Urnenwahl gemietet werden. Abschließend muss
        der Kanzler über den Vorschlag entscheiden.
    \end{itemize}

    Gunhild gibt eine Veränderung der Zusammensetzung des Parlamentes bekannt:
    Daniel Runge ist zurückgetreten. Für ihn rückt Malte Peters (beide BART-Liste)
    nach.

    \subsubsection{Anfragen an das Präsidium}

    Till fragt nach, wie das Prozedere bzgl. der Abholung der Briefwahlstimmen
    aussehe.
    Gunhild erklärt, dass es noch keine konkreten Überlegungen gebe; ein Großteil
    werde jedoch schon in der ersten Januarwoche erledigt.

    \subsubsection{Geschäftsbericht AStA}

    Franziska berichtet von einer Versammlung und Aktion gegen die Schuldenbremse.
    Ferner gibt sie folgende Termine bekannt:
    \begin{itemize}
        \item 05. Mai 2017: Tagung zur Rahmenprüfungsordnung
        \item 17. Mai 2017: Dies Academicus zur Studienreform
        \item 01. November 2017: Dies Academicus zur Hochschulfinanzierung
    \end{itemize}

    Artur erklärt, es habe eine Sitzung der Vertreterversammlung des
    Studierendenwerkes stattgefunden. Insbesondere sei hierbei thematisiert worden,
    dass der Etat des Studierendenwerkes um jährlich 50.000 Euro ansteigen solle.
    Ferner wolle sich das Studierendenwerk kundenorientierter ausrichten. Seitens
    des AStA seien zwei Anträge gestellt worden, zum einen bezüglich des
    Instandhaltungsbedarfes im Rudolf-Laun-Haus, sowie zum anderen hinsichtlich
    der "demokratischen Kultur" in der Vertreterversammlung selbst.

    Philipp berichtet von einem Gespräch mit den Schulleiterinnen und Schulleitern
    der Hamburger Stadtteilschulen bezüglich der Inklusion an ihren Schulen. Der
    AStA erhoffe sich, dass die Aktivitäten der Stadtteilschulleiterinnen und
    -schulleiter mit anderen Aktivitäten zu einem breiten Bildungsbündnis in der
    Stadt zusammengeführt werden können.

    Armin schildert von den Planungen zu den Protestaktionen anlässlich des
    G20-Gipfels. Eine erste Demonstration sei für Februar 2017 geplant.

    \subsubsection{Anfragen an den AStA}

    Keine.

    \subsubsection{Dringlichkeitsanträge des AStA}

    Keine.

    \subsubsection{Feststellung der endgültigen Fassung des Teils B der Tagesordnung}


    %\newpage
    %%% BESCHLOSSENE TAGESORDNUNG %%%
    \underline{Endgültige Fassung der Tagesordnung}
    \begin{enumerate}[label={\textbf{Top \theenumi}},leftmargin=*]
        \item Beitragsordnung (V1617-046) (30 Min.)
        \item VG-Wort (V1617-047) (30 Min.)
        \item Verschiedenes (5 Min.)
    \end{enumerate}

    Die Tagesordnung wird, wie seitens des Präsidiums vorgelegt, einstimmig
    beschlossen.

    \subsubsection{Feststellung der Beschlussfähigkeit}

    Mit 19 anwesenden Mitgliedern ist das Parlament nicht beschlussfähig.

    \textit{Jim beendet die Sitzung um 21.30 Uhr.}
\end{document}
